\pdfvariable suppressoptionalinfo 611
\providecommand{\CRevisionRound}{2}
\documentclass[10pt]{article}
\usepackage[a4paper,margin=17mm]{geometry}
\usepackage{amsmath,amssymb,amsthm,booktabs,microtype}
\usepackage[T1]{fontenc}
\usepackage{lmodern}
\usepackage[hidelinks,hypertexnames=false]{hyperref}
\newtheorem{theorem}{Theorem}
\newtheorem{proposition}[theorem]{Proposition}
\newcommand{\R}{\mathbb R}
\newcommand{\dd}{\,\mathrm d}
\title{A Clean Primitive-Family Atlas for the Circular Billiard}
\author{HCS Research Program}
\date{30 August 2026\quad (revision \CRevisionRound)}
\begin{document}
\maketitle
\begin{abstract}
In the disk, the billiard map is exactly a rigid translation in the boundary
angle and signed half-chord-angle coordinates.  We prove one all-parameter
closure theorem: every reduced fraction \((m/n)\), with
\(1\le m<n/2\), gives two orientation-separated clean \(S^1\) families,
with explicit chord length,
caustic, action, repetitions, and unipotent return shear.  The diameter and
grazing faces are treated separately.  The auxiliary \(p=\sin\alpha\) is not
called a canonical momentum.  This is a geometric Route-A result; no target
arithmetic or spectral match is asserted.
\end{abstract}

\section{Frozen disk and incidence convention}
Let \(D_R=\{x^2+y^2<R^2\}\).  At a collision write
\(\theta\in\R/2\pi\mathbb Z\)
for boundary angle and \(\alpha\in(-\pi/2,\pi/2)\) for the signed half-chord
angle defined by the oriented central increment
\(\theta'-\theta=2\alpha\pmod {2\pi}\).  Its absolute value is the acute
angle-to-tangent magnitude and its sign records direction.  The auxiliary
incidence amplitude
\(p=\sin\alpha\) records the sign and angle magnitude; it is not a
canonical momentum.  In the working angle chart the exact map is
\begin{equation}
 B(\theta,\alpha)=(\theta+2\alpha,\alpha)\pmod {2\pi},
 \qquad DB=\begin{pmatrix}1&2\\0&1\end{pmatrix}.       \label{eq:map}
\end{equation}
The physical billiard symplectic form is expressed using the appropriate
canonical boundary momentum; we make no claim that
\(\dd\theta\wedge\dd\alpha\) is that form.

\ifnum\CRevisionRound=0
The baseline receipt records (\ref{eq:map}) and the fundamental rational
range.  Endpoint faces and clean-return kernels are reserved for the later
rounds.
\fi

\section{Primitive rational families}
\begin{theorem}[complete interior rational families]\label{thm:families}
For \(\gcd(m,n)=1\) and \(1\le m<n/2\), put
\(\alpha_\varepsilon=\varepsilon\pi m/n\),
\(\varepsilon\in\{+1,-1\}\).  Then
\(B^n(\theta,\alpha_\varepsilon)=(\theta,\alpha_\varepsilon)\) for every
\(\theta\), and \(n\) is the minimal positive bounce period.  Conversely,
every interior rational rotation family has this form.
\end{theorem}
\begin{proof}
Iterating \eqref{eq:map} gives
\(B^q(\theta,\alpha)=(\theta+2q\alpha,\alpha)\).
For \(\alpha=\varepsilon\pi m/n\), the shift after \(n\) bounces is
\(\varepsilon2\pi m\).  If a smaller \(q\) closes, \(qm/n\in\mathbb Z\),
which contradicts \(\gcd(m,n)=1\).  Conversely, a rational rotation in the
fundamental interval has a unique reduced numerator/denominator and one of
the two signs.
\end{proof}

\ifnum\CRevisionRound>0
\begin{proposition}[length, caustic, and action]\label{prop:geometry}
For either orientation the chord, polygon length, concentric caustic radius,
and fixed-speed geometric action are
\begin{equation}
 \ell_{m,n}=2R\sin(\pi m/n),\quad
 L_{m,n}=2nR\sin(\pi m/n),\quad
 r_c=R\cos(\pi m/n),\quad S_{m,n}=p_0L_{m,n}.          \label{eq:geometry}
\end{equation}
At the frozen unit speed \(p_0=1\), \(S=L\).  The \(k\)-fold repetition has
\(kn\) bounces and multiplies both \(L\) and \(S\) by \(k\); its unreduced
label \((km,kn)\) is retained but is not primitive.
\end{proposition}
\begin{proof}
The central angle between successive endpoints is \(2\pi m/n\), so the chord
is the corresponding isosceles-triangle base.  Its distance from the center is
\(R\cos(\pi m/n)\), proving the caustic statement.  Summing \(n\) equal
chords and integrating unit-speed momentum gives \eqref{eq:geometry}; repetition
is immediate.
\end{proof}
\fi

\begin{table}[t]
\centering\small
\caption{Selected rows from the finite exact/high-precision receipt (R=1).}
\begin{tabular}{c|c|c|c|c|c}
\toprule
$(m,n)$ & orientation & $\alpha$ & $\ell$ & $L=S$ & $r_c$\\
\midrule
$(1,3)$ & $+$ & $\pi/3$ & $1.7321$ & $5.1962$ & $0.5000$\\
$(1,4)$ & $-$ & $-\pi/4$ & $1.4142$ & $5.6569$ & $0.7071$\\
$(2,5)$ & $+$ & $2\pi/5$ & $1.9021$ & $19.0211$ & $0.3090$\\
$(3,7)$ & $-$ & $-3\pi/7$ & $1.9499$ & $27.2986$ & $0.2225$\\
\bottomrule
\end{tabular}
\end{table}

\section{Clean return and determinant obstruction}
For a primitive row, the fixed set is the one-dimensional family
\(\operatorname{Fix}(B^n)\supset\{(\theta,\alpha_\varepsilon):\theta\in S^1\}\).
In the angle chart,
\begin{equation}
 DB^n=\begin{pmatrix}1&2n\\0&1\end{pmatrix},\qquad
 DB^n-I=\begin{pmatrix}0&2n\\0&0\end{pmatrix}.         \label{eq:shear}
\end{equation}
Hence
\begin{equation}
 \ker(DB^n-I)=\operatorname{span}\{(1,0)\}=T(S^1_\theta),\quad
 \det(I-DB^n)=0.                                      \label{eq:kernel}
\end{equation}
The return is unipotent, with both eigenvalues one.  Thus an isolated-orbit
determinant denominator is obstructed: the clean family must not be assigned
an isolated amplitude.  The finite receipt verifies the kernel itself, not
only the vanishing determinant.

\ifnum\CRevisionRound=1
\section{Revision-one ledger}
The producer enumerates every primitive pair through \(n=12\), both signs,
six explicit repetitions, and the endpoint rows.  Chebyshev relations
\(T_n(\cos(\pi m/n))=(-1)^m\) provide an algebraic/high-precision receipt.
An independent checker makes more than two thousand assertions and a SymPy
script repeats the trigonometric and shear identities.
\fi

\section{Diameter, grazing, and natural quantization}
The diameter case \((m,n)=(1,2)\), \(\alpha=\pm\pi/2\), is one
endpoint-equivalent family, not two families.  Its angle-chart return matrix
for two bounces is the boundary value
\(\left[\begin{smallmatrix}1&4\\0&1\end{smallmatrix}\right]\), but it is not an
interior regular row.  At \(\alpha=0\), \(p=0\), the chord has zero length;
the receipt records one grazing row with two one-sided oriented limits.  No
boundary degeneration is silently counted as a new primitive flight.

For A4 we record only the standard self-adjoint quantizations of the same disk:
\(-\Delta_{D_R}\) with \(f|_{\partial D_R}=0\) (Dirichlet), or with
\(\partial_\nu f|_{\partial D_R}=0\) (Neumann).  This is a definition of a
source-local quantization problem, not an eigenvalue table or a target match.

\ifnum\CRevisionRound>1
\section{Verification and Route-A boundary}
The final receipt contains 44 primitive rows, six repetitions, two boundary
rows, 13 exact identities, independent checks, byte replay, and 31 hostile
repaired-hash mutations.  The frozen scope is
\texttt{NO\_BAD\_EULER\_OR\_ROOT\_NUMBER}; no prime/zero table, Euler factor,
root number, automorphy claim, target determinant, or Hilbert--Pólya operator
is used.  Therefore the locked tuple is
\texttt{(A0\_FAIL,A1\_PASS\_ANALYTIC,A2\_FAIL,A3\_FAIL,
A4\_NATURAL\_QUANTIZATION)},
with overall verdict \texttt{ROUTE\_A\_REJECTED}.
\fi

\begin{thebibliography}{9}
\bibitem{Birkhoff} G.~D. Birkhoff, ``On the periodic motions of dynamical
systems,'' \emph{Acta Mathematica} 50 (1927), 359--379, DOI
\href{https://doi.org/10.1007/BF02421325}{10.1007/BF02421325}.
\bibitem{Bishop} R.~L. Bishop, ``Circular Billiard Tables, Conjugate Loci,
and a Cardioid,'' \emph{Regular and Chaotic Dynamics} 8 (2003), 83--95,
DOI \href{https://doi.org/10.1070/RD2003v008n01ABEH000227}
{10.1070/RD2003v008n01ABEH000227}.
\end{thebibliography}
\end{document}
